\documentclass[10pt]{amsart}
\usepackage{calc,amssymb,amsmath,ulem,stmaryrd,fullpage}
\usepackage{enumerate}
\usepackage{alltt}
\usepackage{multicol}
\RequirePackage[dvipsnames,usenames]{color}
\let\mffont=\sf
\normalem
%\input{kmacros3.sty}
\input{xy}
\xyoption{all}
\usepackage{hyperref}
\usepackage{graphicx}
\usepackage[all,cmtip]{xy}
\usepackage{verbatim}
\usepackage[left=0.95in,top=0.95in,right=0.95in,bottom=0.95in]{geometry}
\newcommand{\tauCohomology}{T}
\newcommand{\FCohomology}{S}


\theoremstyle{theorem}
%\newtheorem*{mainthm*}{Main Theorem}

\newcommand{\note}[1]{\marginpar{\sffamily\tiny #1}}

\def\todo#1{\textcolor{Mahogany}%
{\sffamily\footnotesize\newline{\color{Mahogany}\fbox{\parbox{\textwidth-15pt}{#1}}}\newline}}

\renewcommand{\O}{\mathcal O}
\newcommand{\ie}{{\itshape i.e.} }
\newcommand{\roundup}[1]{\lceil #1 \rceil}
\newcommand{\rounddown}[1]{\lfloor #1 \rfloor}
\renewcommand{\to}[1][]{\xrightarrow{\ #1\ }}
\newcommand{\onto}[1][]{\protect{\xrightarrow{\ #1\ }\hspace{-0.8em}\rightarrow}}
\newcommand{\into}[1][]{\lhook \joinrel \xrightarrow{\ #1\ }}
%\newcommand{DBDual}[1]{{\underline{\omega}_{#1}}}
\newcommand{\DBDual}[1]{{{\uuline \omega}{}^{\mydot}_{#1}}}
\newcommand{\Tt}{{\mathfrak{T}}}
%\newcommand{\bn}{{\mathfrak{n}} }
\newcommand{\sol}[1]{ \vskip 6pt{\color{blue} {\bf Solution:} #1} \vskip 6pt}

\newcommand{\bR}{{\mathbb{R}}}
\newcommand{\va}{{\vec{a}}}
\newcommand{\vb}{{\vec{b}}}
\newcommand{\vzero}{{\vec{0}}}
\newcommand{\vi}{{\vec{i}}}
\newcommand{\vj}{{\vec{j}}}
\newcommand{\vk}{{\vec{k}}}
\newcommand{\vh}{{\vec{h}}}
\newcommand{\vr}{{\vec{r}}}
\newcommand{\vu}{{\vec{u}}}
\newcommand{\vv}{{\vec{v}}}
\newcommand{\vw}{{\vec{w}}}
\newcommand{\vx}{{\vec{x}}}
\newcommand{\vy}{{\vec{y}}}
\newcommand{\vz}{{\vec{z}}}
\newcommand{\vT}{{\vec{T}}}
\newcommand{\vN}{{\vec{N}}}
\newcommand{\vF}{{\vec{F}}}
\newcommand{\vG}{{\vec{G}}}
\newcommand{\bF}{\mathbb{F}}
\newcommand{\bQ}{\mathbb{Q}}
\newcommand{\bZ}{\mathbb{Z}}
\newcommand{\bC}{\mathbb{C}}
\newcommand{\Aut}{\mathrm{Aut}}


\begin{document}
\title{(Review) Worksheet \#6-- Math 5320\\Spring 2020}
\author{Not due}
\maketitle

I'm trying to make these questions like those I might put on the final.
\noindent
{\bf 1.}  Short answer.
\vskip 3pt
{\bf (a)}  Give an example of a ring that is not an integral domain.
\vskip 75pt
{\bf (b)}  Give an example of a prime ideal that is not maximal.
\vskip 75pt
{\bf (c)}  Suppose that $R$ is an integral domain.  What does it mean for $r \in R$ to be irreducible? 
\vskip 75pt
{\bf (d)}  Carefully state Eisenstein's irreducibility criterion for polynomials in $\bQ[x]$.
\vskip 75pt
{\bf (e)}  Suppose $f(x) \in \bZ[x]$ is a monic polynomial and it is irreducible.  Let $f_7(x) \in \bZ/7\bZ[x]$ by its reduction to characteristic $7$.  Is $f_7(x)$ always irreducible?
\vskip 75pt
{\bf (f)}  Suppose that $R$ is an integral domain and $r \in R$ is a prime element.  Prove that $r$ is irreducible.
\newpage
\noindent
{\bf 2.}  More short answer.
\vskip 3pt
{\bf (a)}  Give an example of an extension of fields $\bQ \subseteq K$ such that $[K : \bQ] = 5$.
\vskip 75pt
{\bf (b)}  What does it mean that a finite extension $F \subseteq K$ of fields in characteristic zero is \emph{Galois}?
\vskip 75pt
{\bf (c)}  Is there a field with $51$ elements?
\vskip 75pt
{\bf (d)}  Suppose $p$ is a prime number and $e > 0$.  Write down a polynomial such that every element of $\bF_{p^e}$ is a root of that polynomial.
\vskip 75pt
{\bf (e)}  Suppose that $\bF_9$ is the field with 9 elements.  What finite field $K$ has the property that $[K : \bF_9] = 3$?
\vskip 75pt
{\bf (f)}  Give an example of a field extension $F \subseteq K$ such that $\Aut(K/F)$ is not Abelian.
\newpage
{\bf 3.}  Even more short answer.
\noindent
\vskip 3pt
{\bf (a)}  Is $\bZ[x]$ a PID?
\vskip 75pt
{\bf (b)}  If $\phi : R \to S$ is a ring homomorphism and $I \subseteq S$ is an ideal, is it always true that $\phi^{-1}(I)$ is an ideal?
\vskip 75pt
{\bf (c)}  How many elements are in the ring $\bZ/6\bZ[x]/\langle x^2 + 1 \rangle$?
\vskip 75pt
{\bf (d)}  Give an example of a ring that is not a UFD.
\vskip 75pt
{\bf (e)}  List at least 4 different numbers in $\bR$ that are \emph{not} constructible over $\bQ$ but are algebraic over $\bQ$.
\vskip 75pt
{\bf (f)}  State the Main Theorem of Galois Theory.
\newpage
\noindent
{\bf 4.}  Suppose $\phi : R \to S$ is a ring homomorphism and $P \subseteq S$ is a prime ideal.  Prove that $\phi^{-1}(P)$ is a prime ideal of $R$.
\vskip 275pt
{\bf 5.}  Consider the finite field $K = \bF_2[x]/\langle x^2 + x + 1 \rangle$.  Write down addition and multiplication tables for $K$.  
\newpage
\noindent
{\bf 6.}  Suppose that $k$ is a field.  Prove that $k[x]$ is a PID.
\vskip 275pt
\noindent
{\bf 7.}  Suppose $F \subseteq K$ is an extension of fields and $H \subseteq \Aut(K/F)$ is a subgroup.  Carefully write down a proof that $K^H$ is a field.  
\newpage
\noindent
{\bf 8.}  Let $K$ be the splitting field of $x^4 + 3$ over $\bQ$. Compute $[K : \bQ]$.  Find generators for the group $\Aut(K/F)$, is this group Abelian?  
\vskip 275pt
\noindent
{\bf 9.}  Find 2 proper nontrivial subgroups of $\Aut(K/\bQ)$ (with $K$ as above) and compute their associated fixed fields. 
\newpage
{\bf 10.}  How many irreducible degree 4 polynomials are there over $\bF_5$.  
\vskip 275pt
{\bf 11.}  Let $R$ be an integral domain and suppose that $W$ is a multiplicative set.  Suppose that $P \subseteq R$ is a prime ideal and that $P \cap W = \emptyset$.  Prove that $W^{-1}P$ is not equal to $W^{-1}R$.
\newpage
\noindent
{\bf 12.}  Suppose $p$ is prime and $e > 0$.  Give a careful proof that $\bF_{p^e}^{\times} = \bF_{p^e} \setminus \{0\}$ is a cyclic group.
\vskip 275pt
\noindent
{\bf 13.}  Which maximal ideals in the polynomial ring $\bC[x,y]$ contain both $x^2 + y^2 -4$ and $xy - 1$?
\end{document} 